\chapter{Dados em Painel}
\label{cap:painel}
% ── índice ──────────────────────────────────────────────────────
\index{dados em painel|textbf}
\index{efeitos fixos|textbf}
\index{efeitos aleatórios|textbf}
\index{within estimator|see{efeitos fixos}}
\index{fe@\texttt{fe()}|textbf}
\index{re@\texttt{re()}|textbf}
\index{Hausman, teste de|textbf}
\index{hausman@\texttt{hausman()}|textbf}
\index{viés de Nickell}

% ─────────────────────────────────────────────────────────────────────────────
\section{O modelo}

Um painel observa $n$ indivíduos ao longo de $T$ períodos. O modelo geral é:
\[
  y_{it} = \alpha_i + X_{it}\beta + \varepsilon_{it},
  \qquad i = 1,\ldots,n;\quad t = 1,\ldots,T
\]

onde $\alpha_i$ é o \textbf{efeito individual não observado} — características
fixas no tempo que diferem entre indivíduos (habilidade, cultura, localização).
O OLS em corte transversal ou empilhado ignora $\alpha_i$, tratando-o como parte
do erro:

\[
  \mathrm{plim}\;\hat\beta_{\mathrm{OLS}}
  = \beta + \frac{\mathrm{Cov}(x_{it},\,\alpha_i)}{\mathrm{Var}(x_{it})}
\]

O viés só desaparece se $\mathrm{Cov}(x_{it}, \alpha_i) = 0$ — hipótese forte
na maioria das aplicações.

% ─────────────────────────────────────────────────────────────────────────────
\section{Efeitos Fixos}

O estimador \textbf{within} elimina $\alpha_i$ subtraindo a média temporal
de cada indivíduo:
\[
  \tilde{y}_{it} = y_{it} - \bar{y}_i, \qquad
  \tilde{X}_{it} = X_{it} - \bar{X}_i
\]

O estimador de Efeitos Fixos é então o OLS sobre as variáveis desmediadas:
\[
  \boxed{\hat\beta_{\mathrm{FE}} = (\tilde{X}'\tilde{X})^{-1}\tilde{X}'\tilde{y}}
\]

Como a transformação within anula $\alpha_i$ para qualquer valor de
$\mathrm{Cov}(x_{it}, \alpha_i)$, o estimador é consistente mesmo sob
endogeneidade do efeito individual. A variância é:
\[
  \widehat{\mathrm{Var}}[\hat\beta_{\mathrm{FE}}] = s^2_{\mathrm{FE}}(\tilde{X}'\tilde{X})^{-1},
  \qquad
  s^2_{\mathrm{FE}} = \frac{\hat\varepsilon'\hat\varepsilon}{nT - n - k}
\]

\begin{aviso}
  Variáveis que não variam no tempo (sexo, raça, região de nascimento) são
  eliminadas pela transformação within. FE não pode estimar seu efeito.
  Use RE ou estimadores híbridos nesses casos.
\end{aviso}

% ─────────────────────────────────────────────────────────────────────────────
\section{Efeitos Aleatórios}

Se $\alpha_i \sim \mathrm{IID}(0, \sigma^2_\alpha)$ e
$\mathrm{Cov}(x_{it}, \alpha_i) = 0$, $\alpha_i$ é parte do erro composto
$v_{it} = \alpha_i + \varepsilon_{it}$. A estrutura de erros implica GLS com
quasi-desmediação:
\[
  y^*_{it} = y_{it} - \theta\bar{y}_i,
  \qquad
  \theta = 1 - \frac{\sigma_\varepsilon}{\sqrt{\sigma^2_\varepsilon + T\sigma^2_\alpha}}
  \in [0,1)
\]

O estimador RE é o OLS sobre $y^*_{it}$ e $X^*_{it}$. É mais eficiente que
FE sob $H_0: \mathrm{Cov}(x_{it},\alpha_i) = 0$, mas \emph{inconsistente}
quando a hipótese falha.

% ─────────────────────────────────────────────────────────────────────────────
\section{Teste de Hausman}

O teste de Hausman compara FE (consistente sempre) com RE (consistente sob
$H_0$, eficiente):
\[
  H_0: \mathrm{Cov}(x_{it}, \alpha_i) = 0
  \quad\Longrightarrow\quad
  \mathrm{plim}(\hat\beta_{\mathrm{FE}} - \hat\beta_{\mathrm{RE}}) = 0
\]

A estatística é:
\[
  H = (\hat\beta_{\mathrm{FE}} - \hat\beta_{\mathrm{RE}})'
      \bigl[\widehat{\mathrm{Var}}(\hat\beta_{\mathrm{FE}}) -
             \widehat{\mathrm{Var}}(\hat\beta_{\mathrm{RE}})\bigr]^{-1}
      (\hat\beta_{\mathrm{FE}} - \hat\beta_{\mathrm{RE}})
  \;\xrightarrow{d}\; \chi^2(k)
\]

Rejeitar $H_0$ indica que RE é inconsistente — use FE.

% ─────────────────────────────────────────────────────────────────────────────
\section{Verificação por DGP}

\subsection*{Recuperação exata (sem ruído)}

O DGP usa 100 indivíduos × 10 períodos ($n = 1000$). O efeito fixo
$\alpha_i = 0{,}05\cdot i$ é determinístico e correlacionado com $x$:
\[
  x_{it} = \alpha_i + u_{it},\quad u \sim \mathcal{N}(0,1)
  \qquad
  y_{it} = 1{,}0 + 2{,}0\,x_{it} + \alpha_i \quad(\varepsilon = 0)
\]

Sem ruído idiossincrático a transformação within elimina $\alpha_i$
exatamente, e FE recupera $\beta$ com precisão de máquina:

\begin{lstlisting}[caption={FE com efeito fixo correlacionado — sem ruído}]
seed(42)
input df
id
1
end
for i in 1..=10 { df = append(df, df) }
df |> filter(_n <= 1000)
generate df id    = (_n - 1) % 100 + 1
generate df alpha = id * 0.05
generate df u     = rnormal()
generate df x     = alpha + u
generate df y     = 1.0 + 2.0*x + alpha
let m_ols = ols(y ~ x, df)
let m_fe  = fe(y ~ x, df, id=id)
esttab(m_ols, m_fe)
\end{lstlisting}

\begin{lstlisting}[language={}, basicstyle=\ttfamily\small,
  frame=none, numbers=none, backgroundcolor=\color{codbg},
  xleftmargin=0pt, xrightmargin=0pt, breaklines=false]
                            (1)              (2)
                            OLS               FE
────────────────────────────────────────────────
x                     2.9663***        2.0000***
                       (0.0060)         (0.0000)
Constant              0.6002***
                       (0.0202)
────────────────────────────────────────────────
N                          1000             1000
R²                       0.9959
Adj. R²                  0.9959
────────────────────────────────────────────────
* p<0.10  ** p<0.05  *** p<0.01
\end{lstlisting}

O OLS superestima $\beta$ em $+0{,}97$ porque $\mathrm{Cov}(x,\alpha) =
\mathrm{Var}(\alpha) > 0$. O FE recupera $2{,}0000$ exatamente.

\subsection*{Hausman com ruído}

Com ruído idiossincrático $\varepsilon \sim \mathcal{N}(0,1)$, o FE permanece
consistente e o RE é inconsistente (pois $\alpha_i$ ainda é correlacionado
com $x$). O Hausman deve rejeitar $H_0$:

\begin{lstlisting}[caption={Hausman — FE vs RE com efeito correlacionado}]
// mesmo DGP; adiciona ruído idiossincrático
generate df e  = rnormal()
generate df y2 = 1.0 + 2.0*x + alpha + e
let m_fe2 = fe(y2 ~ x, df, id=id)
let m_re2 = re(y2 ~ x, df, id=id)
esttab(m_fe2, m_re2)
hausman(m_fe2, m_re2)
\end{lstlisting}

\begin{lstlisting}[language={}, basicstyle=\ttfamily\small,
  frame=none, numbers=none, backgroundcolor=\color{codbg},
  xleftmargin=0pt, xrightmargin=0pt, breaklines=false]
                            (1)              (2)
                             FE               RE
────────────────────────────────────────────────
x                     2.0977***        2.9347***
                       (0.0338)         (0.0118)
const                                  1.1930***
                                        (0.0399)
────────────────────────────────────────────────
N                          1000                0
────────────────────────────────────────────────
* p<0.10  ** p<0.05  *** p<0.01

══════════════════════════════════════════════════════════════
 Hausman Test: FE vs RE
 H₀: efeitos individuais não correlacionados com regressores
══════════════════════════════════════════════════════════════
── Coeficientes Comuns
   Variável         β_FE       β_RE         Δβ
   x              2.0977     2.9347    -0.8369
── Estatística
   χ²(1) = 698.1711   p = 0.0000  ***
── Conclusão
   Rejeita H₀ → use EFEITOS FIXOS (RE pode ser inconsistente)
══════════════════════════════════════════════════════════════
\end{lstlisting}

O RE absorve o efeito de $\alpha_i$ no coeficiente de $x$ (2.9347), replicando
o viés do OLS. O Hausman rejeita $H_0$ com $p < 0{,}001$ — a endogeneidade
do efeito individual é detectada.\footnote{%
  O FE recupera 2.0977, não 2.0000, porque a transformação within elimina
  $\alpha_i$ mas não $\varepsilon_{it}$. Com ruído idiossincrático presente,
  o estimador é consistente mas não exato — o mesmo motivo pelo qual OLS
  e IV divergem de seus parâmetros verdadeiros quando $\varepsilon \neq 0$.}

% ─────────────────────────────────────────────────────────────────────────────
\section{Sintaxe}

\begin{lstlisting}
fe(fórmula, df, id=col)
re(fórmula, df, id=col)
hausman(modelo_fe, modelo_re)
\end{lstlisting}

A coluna \lstinline{id} identifica o indivíduo. O FE remove automaticamente
o intercepto (absorvido pelos efeitos individuais).

\begin{lstlisting}[caption={Fluxo completo}]
load "painel.csv" as df

let m_fe = fe(salario ~ educ + exp, df, id=id)
let m_re = re(salario ~ educ + exp, df, id=id)

esttab(m_fe, m_re)
hausman(m_fe, m_re)
\end{lstlisting}

\section{Opções de erro padrão}

\begin{lstlisting}
fe(y ~ x, df, id=id)                        // clássico
fe(y ~ x, df, id=id, cov=robust)            // HC1
// cluster por indivíduo
fe(y ~ x, df, id=id, cov=cluster, cluster=id)
\end{lstlisting}

\begin{nota}
  Em painéis, erros clusterizados por indivíduo (\lstinline{cluster=id})
  são o padrão recomendado: corrigem heterocedasticidade e correlação serial
  dentro do indivíduo simultaneamente.
\end{nota}

\section{Capturando resultados}

\begin{lstlisting}
let m = fe(y ~ x, df, id=id)
display m
let yhat = predict(m, "xb")
\end{lstlisting}

\section{Diagnósticos de painel}

A escolha entre OLS agrupado, Efeitos Fixos e Efeitos Aleatórios depende de
testes formais. Hayashi implementa todos os testes padrão da literatura.

\subsection{OLS agrupado vs.\ Efeitos Aleatórios}

\textbf{Breusch-Pagan LM} (\texttt{bptest}): testa $H_0: \sigma_u^2 = 0$
(não há efeito de painel — OLS agrupado é adequado):

\begin{lstlisting}
bptest(df, y ~ x, id="entity")
\end{lstlisting}

Se rejeitar $H_0$, use RE ou FE em vez de OLS agrupado.

\subsection{OLS agrupado vs.\ Efeitos Fixos}

\textbf{F-test para FE} (\texttt{ftest\_fe}): testa $H_0$: todos os efeitos
individuais são zero:

\begin{lstlisting}
ftest_fe(df, y ~ x, id="entity")
\end{lstlisting}

Rejeitar $H_0$ indica que FE é preferível a OLS agrupado.

\subsection{Efeitos Fixos vs.\ Efeitos Aleatórios}

\textbf{Hausman} (\texttt{hausman}): testa $H_0$: RE é consistente (efeitos
individuais não correlacionados com regressores):

\begin{lstlisting}
hausman(m_fe, m_re)
\end{lstlisting}

Rejeitar $H_0$ implica FE (consistente sob correlação); não rejeitar implica
RE (mais eficiente).

\textbf{Mundlak} (\texttt{mundlak}): generalização do Hausman que aumenta o
modelo com médias individuais dos regressores:

\begin{lstlisting}
mundlak(df, y ~ x1 + x2, id="entity")
\end{lstlisting}

\subsection{Correlação serial e dependência cross-section}

\textbf{Wooldridge} (\texttt{wooldridge}): testa $H_0$: ausência de
correlação serial de primeira ordem nos erros idiossincráticos:

\begin{lstlisting}
wooldridge(df, y ~ x, id="entity", time="time")
\end{lstlisting}

\textbf{Pesaran CD} (\texttt{pesaran}): testa $H_0$: ausência de dependência
cross-sectional. Relevante para painéis com $N$ grande onde spillovers
espaciais ou de rede podem existir:

\begin{lstlisting}
pesaran(df, y ~ x, id="entity", time="time")
\end{lstlisting}

\subsection{Arellano-Bond m1/m2}

\textbf{Arellano-Bond} (\texttt{abtest}): testa correlação serial nos
resíduos em primeira diferença. $m_1$ deve rejeitar $H_0$ (a FD induz AR(1)
por construção); $m_2$ não deve rejeitar $H_0$ (valida os instrumentos
$y_{i,t-2}$ do GMM):

\begin{lstlisting}
abtest(df, y ~ x, id="entity", time="time")
\end{lstlisting}
